
\section{\label{sec:Speed}Kecepatan}

Selain akurasi, tidak ada hal yang lebih penting dalam metode numerik
daripada kecepatan. Namun, hampir selalu terdapat kompromi di antara
keduanya. Perhitungan cepat sering kali tidak terlalu akurat, sedangkan
perhitungan akurat sering kali tidak terlalu cepat. Meskipun demikian,
algoritme tertentu menghasilkan nilai akurat dengan cepat. Menurunkan
algoritme semacam itu, atau mengenalinya setelah diturunkan, merupakan
inti analisis numerik.

Jenis metode numerik pertama yang akan kita temui menghasilkan barisan
hampiran yang, jika semuanya bekerja dengan baik, mendekati suatu
nilai yang diinginkan, misalnya $p$. Dengan metode ini, kita akan
memperoleh barisan $\langle p_{n}\rangle$ dengan $\lim_{n\to\infty}p_{n}=p$.
Anda seharusnya sudah mengenal konsep limit barisan dari Kalkulus,
tetapi tujuannya di sana sangat berbeda dari tujuan kita di sini.
Umumnya, perhatian Anda tertuju pada apakah suatu barisan konvergen
atau tidak. Jika barisan tersebut konvergen dan Anda sangat beruntung,
Anda dapat menentukan limitnya. Dalam analisis numerik, kita mengetahui
bahwa barisan tertentu konvergen dan hanya tertarik pada seberapa
cepat kekonvergenan itu terjadi. 

Pengamatan sederhana (dan sedikit akal sehat) dapat menunjukkan mobil
mana di jalan raya yang melaju lebih cepat daripada mobil lainnya.
Pengamatan sederhana (dan sedikit akal sehat) juga sering dapat menunjukkan
barisan mana yang konvergen lebih cepat daripada barisan lainnya.
\begin{table}[h]
\hrule \caption{\label{tab:convergenceTo2.718281828...}Beberapa barisan yang konvergen
ke $e$.}
\medskip{}

\centering{}%
\begin{tabular}{c|cccc}
$n$ & $q_{n}$ & $r_{n}$ & $s_{n}$ & $t_{n}$\tabularnewline
\hline 
0 & 3 & 3 & 3 & 3\tabularnewline
1 & $2.9436563656918$ & $2.86799618929986$ & $2.82129001274358$ & $2.78177393100014$\tabularnewline
2 & $2.89858145824525$ & $2.78315514435127$ & $2.73850656616954$ & $2.72150682612711$\tabularnewline
3 & $2.86252153228801$ & $2.73974041668143$ & $2.71973377603211$ & $2.71829014894701$\tabularnewline
4 & $2.83367359152222$ & $2.72324781752852$ & $2.71830229432561$ & $2.71828182851442$\tabularnewline
5 & $2.81059523890958$ & $2.71899828870116$ & $2.71828184916891$ & $2.71828182845904$\tabularnewline
6 & $2.79213255681947$ & $2.71833715075158$ & $2.71828182845934$ & $2.71828182845904$\tabularnewline
7 & $2.77736241114739$ & $2.71828369688657$ & $2.71828182845904$ & $2.71828182845904$\tabularnewline
8 & $2.76554629460972$ & $2.71828184959225$ & $2.71828182845904$ & $2.71828182845904$\tabularnewline
9 & $2.75609340137958$ & $2.71828182851528$ & $2.71828182845904$ & $2.71828182845904$\tabularnewline
10 & $2.74853108679547$ & $2.71828182845907$ & $2.71828182845904$ & $2.71828182845904$\tabularnewline
$\vdots$ & $\vdots$ & $\vdots$ & $\vdots$ & $\vdots$\tabularnewline
\end{tabular}\medskip{}
\hrule 
\end{table}
 Perhatikan barisan-barisan dalam Tabel \ref{tab:convergenceTo2.718281828...}
yang semuanya konvergen ke $e\approx2.71828182845904$. $\langle t_{n}\rangle$
sudah akurat hingga 15 angka signifikan pada suku keenam; $\langle s_{n}\rangle$
sudah akurat hingga 15 angka signifikan pada suku kedelapan; $\langle r_{n}\rangle$
masih belum akurat hingga 15 angka signifikan pada suku kesebelas,
tetapi tampaknya akan mencapai akurasi 15 angka signifikan pada suku
kedua belas; sedangkan $\langle q_{n}\rangle$ baru akurat hingga
2 angka signifikan pada suku kesebelas, sehingga tampaknya memerlukan
jauh lebih dari dua belas suku untuk mencapai akurasi 15 angka signifikan.
Karena semuanya dimulai dari 3, cukup masuk akal untuk mengatakan
bahwa, jika diurutkan dari yang tercepat hingga yang paling lambat,
urutannya adalah $\langle t_{n}\rangle$, $\langle s_{n}\rangle$,
$\langle r_{n}\rangle$, $\langle q_{n}\rangle$. Dan urutan itu memang
benar, sebagaimana akan segera kita lihat. Namun, mengetahui mobil
mana yang lebih cepat berbeda dengan mengetahui kecepatan masing-masing
mobil. Demikian pula, mengetahui barisan mana yang konvergen lebih
cepat berbeda dengan mengetahui seberapa cepat masing-masing barisan
itu konvergen. Untuk mengukur kecepatan mobil tertentu, Anda memerlukan
speedometer atau alat pengukur kecepatan radar. Untuk mengukur orde
kekonvergenan (kecepatan) suatu barisan, Anda memerlukan sebuah definisi
dan sedikit aljabar. 

\subsubsection{Orde kekonvergenan suatu barisan\index{kekonvergenan!orde}}

Misalkan barisan $\langle p_{n}\rangle$ konvergen ke $p$. Kita mengatakan
bahwa $\langle p_{n}\rangle$ konvergen dengan orde $\alpha\geq1$
jika
\[
\lim_{n\to\infty}\frac{\left|p_{n+1}-p\right|}{\left|p_{n}-p\right|^{\alpha}}=\lambda
\]
untuk suatu bilangan real $\lambda>0$. 

Mari kita lihat cara menggunakan definisi ini untuk menghitung orde
kekonvergenan barisan-barisan dalam Tabel \ref{tab:convergenceTo2.718281828...}.
Menurut definisi tersebut, $\alpha$, jika ada, menyatakan kecepatan
(atau orde) kekonvergenan \index{kekonvergenan!orde} suatu barisan.
Sekarang, dengan mengasumsikan bahwa $\alpha$ memang ada, kita memperoleh
$\lim_{n\to\infty}\frac{\left|p_{n+1}-p\right|}{\left|p_{n}-p\right|^{\alpha}}=\lambda$,
sehingga untuk $n$ yang cukup besar,
\[
\frac{\left|p_{n+1}-p\right|}{\left|p_{n}-p\right|^{\alpha}}\approx\frac{\left|p_{n+2}-p\right|}{\left|p_{n+1}-p\right|^{\alpha}}\approx\lambda.
\]
Secara khusus, kita dapat menyelesaikan persamaan tersebut terhadap
$\alpha$ dan memperoleh ${\displaystyle \alpha\approx\frac{\ln\left|\frac{p_{n+2}-p}{p_{n+1}-p}\right|}{\ln\left|\frac{p_{n+1}-p}{p_{n}-p}\right|}}$.

\begin{digression}{Orde kekonvergenan kurang dari atau sama dengan 1?} 

\noindent Tidak ada orde kekonvergenan \index{kekonvergenan!orde}
yang kurang dari satu, sebab jika $\lim_{n\to\infty}\frac{\left|p_{n+1}-p\right|}{\left|p_{n}-p\right|^{\alpha}}=\lambda$
untuk suatu $0<\alpha<1$, maka 
\begin{eqnarray*}
\lim_{n\to\infty}\frac{|p_{n+1}-p|}{|p_{n}-p|} & = & \lim_{n\to\infty}\frac{|p_{n+1}-p|}{|p_{n}-p|^{\alpha}}\cdot|p_{n}-p|^{\alpha-1},
\end{eqnarray*}
yang merupakan kontradiksi. Di satu sisi, uji rasio menyiratkan bahwa
$\lim_{n\to\infty}\frac{|p_{n+1}-p|}{|p_{n}-p|}$ ada dan kurang dari
atau sama dengan 1. Di sisi lain, $\alpha<1\implies\alpha-1<0$ sehingga
ketika $|p_{n}-p|$ bernilai kecil, $|p_{n}-p|^{\alpha-1}$ bernilai
besar. Oleh karena itu, $\lim_{n\to\infty}\frac{|p_{n+1}-p|}{|p_{n}-p|^{\alpha}}\cdot|p_{n}-p|^{\alpha-1}$
tidak ada. Untuk membuktikannya secara ketat, misalkan $M$ adalah
sebarang bilangan real. Maka terdapat $N_{1}$ sedemikian sehingga
$n>N_{1}$ menyiratkan $\frac{|p_{n+1}-p|}{|p_{n}-p|^{\alpha}}>0.9\lambda$.
Terdapat pula $N_{2}$ sedemikian sehingga $n>N_{2}$ menyiratkan
${\displaystyle |p_{n}-p|<\left(\frac{0.9\lambda}{M}\right)^{\frac{1}{1-\alpha}}}$,
sehingga $|p_{n}-p|^{\alpha-1}>\frac{M}{0.9\lambda}$. Dengan mengambil
$N=\max\{N_{1},N_{2}\}$, kita memperoleh bahwa $n>N$ menyiratkan
sekaligus $\frac{|p_{n+1}-p|}{|p_{n}-p|^{\alpha}}>0.9\lambda$ dan
$|p_{n}-p|^{\alpha-1}>\frac{M}{0.9\lambda}$. Oleh karena itu, untuk
$n>N$, berlaku 
\begin{eqnarray*}
\frac{|p_{n+1}-p|}{|p_{n}-p|} & = & \frac{|p_{n+1}-p|}{|p_{n}-p|^{\alpha}}\cdot|p_{n}-p|^{\alpha-1}>0.9\lambda\cdot\frac{M}{0.9\lambda}=M.
\end{eqnarray*}
Dengan demikian, $\lim_{n\to\infty}\frac{|p_{n+1}-p|}{|p_{n}-p|}$
tidak ada. Ketika $\alpha=1$, haruslah $\lambda\leq1$ karena jika
tidak, uji rasio menyiratkan bahwa $\langle|p_{n}-p|\rangle$ divergen
dan, dengan demikian, $\langle p_{n}\rangle$ divergen. 

\noindent\end{digression} 

Sebagai contoh, ${\displaystyle \frac{\ln\left|\frac{q_{2}-e}{q_{1}-e}\right|}{\ln\left|\frac{q_{1}-e}{q_{0}-e}\right|}\approx\frac{\ln\left|\frac{2.8985-e}{2.9436-e}\right|}{\ln\left|\frac{2.9436-e}{3-e}\right|}}\approx1$
dan ${\displaystyle \frac{\ln\left|\frac{q_{10}-e}{q_{9}-e}\right|}{\ln\left|\frac{q_{9}-e}{q_{8}-e}\right|}=\frac{\ln\left|\frac{2.7485-e}{2.7560-e}\right|}{\ln\left|\frac{2.7560-e}{2.7655-e}\right|}}\approx1$.
Jika kita mencoba kelompok lain yang terdiri atas tiga suku berurutan
dari $\langle q_{n}\rangle$, kita memperoleh hasil yang sama. Orde
kekonvergenan\index{kekonvergenan!orde} dari $\langle q_{n}\rangle$
kira-kira 1. Tentu saja, kita memerlukan rumus untuk $|q_{n}-e|$
guna menentukan apakah limit tersebut benar-benar 1, tetapi kita sudah
memiliki beberapa bukti pendukung. Dengan mengulangi perhitungan untuk
$\langle r_{n}\rangle$, $\langle s_{n}\rangle$, dan $\langle t_{n}\rangle$,
kita memperoleh perkiraan orde kekonvergenan masing-masing sebesar
$1.322$, $1.618$, dan $2$. Sekali lagi kita melihat bahwa, jika
diurutkan dari yang tercepat hingga yang paling lambat, urutannya
adalah $\langle t_{n}\rangle$, $\langle s_{n}\rangle$, $\langle r_{n}\rangle$,
$\langle q_{n}\rangle$. 

Jika Anda mencoba menghitung sendiri orde kekonvergenan tersebut,
mungkin Anda menyadari bahwa diperlukan informasi tambahan untuk menggunakan
$s_{n}$ ketika $n>6$ atau $t_{n}$ ketika $n>4$. Semua suku tersebut
bernilai sama di dalam tabel, sehingga rumus untuk $\alpha$ gagal
menghasilkan bilangan real! Tabel yang lebih berguna untuk menghitung
orde kekonvergenan adalah tabel yang memuat galat absolut: 
\begin{table}[h]
\hrule \caption{\label{tab:absoluteErrors}Galat absolut.}
\medskip{}

\centering{}%
\begin{tabular}{c|cccc}
$n$ & $|q_{n}-e|$ & $|r_{n}-e|$ & $|s_{n}-e|$ & $|t_{n}-e|$\tabularnewline
\hline 
0 & $2.817(10)^{-1}$ & $2.817(10)^{-1}$ & $2.817(10)^{-1}$ & $2.817(10)^{-1}$\tabularnewline
1 & $2.253(10)^{-1}$ & $1.497(10)^{-1}$ & $1.03(10)^{-1}$ & $6.349(10)^{-2}$\tabularnewline
2 & $1.802(10)^{-1}$ & $6.487(10)^{-2}$ & $2.022(10)^{-2}$ & $3.224(10)^{-3}$\tabularnewline
3 & $1.442(10)^{-1}$ & $2.145(10)^{-2}$ & $1.451(10)^{-3}$ & $8.32(10)^{-6}$\tabularnewline
4 & $1.153(10)^{-1}$ & $4.965(10)^{-3}$ & $2.046(10)^{-5}$ & $5.538(10)^{-11}$\tabularnewline
5 & $9.231(10)^{-2}$ & $7.164(10)^{-4}$ & $2.07(10)^{-8}$ & $2.453(10)^{-21}$\tabularnewline
6 & $7.385(10)^{-2}$ & $5.532(10)^{-5}$ & $2.953(10)^{-13}$ & $4.817(10)^{-42}$\tabularnewline
7 & $5.908(10)^{-2}$ & $1.868(10)^{-6}$ & $4.263(10)^{-21}$ & $1.856(10)^{-83}$\tabularnewline
8 & $4.726(10)^{-2}$ & $2.113(10)^{-8}$ & $8.777(10)^{-34}$ & $2.757(10)^{-166}$\tabularnewline
9 & $3.781(10)^{-2}$ & $5.623(10)^{-11}$ & $2.608(10)^{-54}$ & $6.084(10)^{-332}$\tabularnewline
10 & $3.024(10)^{-2}$ & $2.22(10)^{-14}$ & $1.595(10)^{-87}$ & $2.961(10)^{-663}$\tabularnewline
\end{tabular}\medskip{}
\hrule 
\end{table}
 Selain mempermudah penghitungan $\alpha$, tabel ini menunjukkan
dengan sangat jelas bahwa simpulan berdasarkan intuisi kita tentang
barisan mana yang berkonvergensi lebih cepat memang benar. Bandingkan
saja akurasi (galat absolut) suku kesebelasnya. 

Sekarang kita dapat menghitung orde kekonvergenan\index{kekonvergenan!orde},
tetapi apa sebenarnya maknanya? Apa yang diberitahukan orde kekonvergenan
tentang suku-suku berurutan dalam barisan? Dengan menyelesaikan hampiran
$\frac{\left|p_{n+1}-p\right|}{\left|p_{n}-p\right|^{\alpha}}\approx\lambda$
kita peroleh $|p_{n+1}-p|\approx\lambda|p_{n}-p|^{\alpha}$. Jadi,
secara kasar, kekonvergenan berorde $\alpha$ berarti bahwa, untuk
$n$ yang cukup besar, hampiran $p_{n+1}$ sekitar $\lambda|p_{n}-p|^{\alpha-1}$
kali lebih dekat ke limit $p$ daripada $p_{n}$. Untuk menyatakannya
kembali dalam banyaknya angka signifikan yang akurat, kita lakukan
sedikit aljabar:
\begin{eqnarray*}
|p_{n+1}-p| & \approx & \lambda|p_{n}-p|^{\alpha}\\
\left|\frac{p_{n+1}-p}{p}\right| & \approx & \lambda\left|\frac{p_{n}-p}{p}\right|^{\alpha}\cdot|p|^{\alpha-1}\\
-\log\left|\frac{p_{n+1}-p}{p}\right| & \approx & -\log\left|\frac{p_{n}-p}{p}\right|^{\alpha}-\log\left(\lambda|p|^{\alpha-1}\right)\\
d(p_{n+1}) & \approx & \alpha d(p_{n})-\log\left(\lambda|p|^{\alpha-1}\right).
\end{eqnarray*}
Berdasarkan perhitungan ini, kita memperoleh aturan praktis berikut:
\begin{enumerate}
\item untuk kekonvergenan linear ($\alpha=1$), $d(p_{n+1})\approx d(p_{n})-\log\lambda$,
sehingga setiap suku memiliki sejumlah tetap angka signifikan yang
akurat lebih banyak (kira-kira sama dengan $-\log\lambda$) dibanding
suku sebelumnya;
\item untuk kekonvergenan kuadratik ($\alpha=2$), $d(p_{n+1})\approx2d(p_{n})-\log\left(\lambda|p|\right)$,
sehingga setiap suku memiliki kira-kira dua kali banyaknya angka signifikan
yang akurat dibanding suku sebelumnya, dengan sedikit penyimpangan;
\item untuk kekonvergenan kubik ($\alpha=3$), $d(p_{n+1})\approx3d(p_{n})-\log\left(\lambda|p|^{2}\right)$,
sehingga setiap suku memiliki kira-kira tiga kali banyaknya angka
signifikan yang akurat dibanding suku sebelumnya, dengan sedikit penyimpangan;
\end{enumerate}
dan seterusnya. Ringkasnya, untuk $n$ yang besar, setiap suku diperkirakan
memiliki $-\log\left(\lambda|p|^{\alpha-1}\right)$ angka signifikan
akurat lebih banyak daripada $\alpha$ kali banyaknya angka signifikan
akurat pada suku sebelumnya. Kita dapat melihat klaim ini secara langsung
dengan menghitung $\lambda$ untuk barisan $\langle t_{n}\rangle$,
$\langle s_{n}\rangle$, $\langle r_{n}\rangle$, dan $\langle q_{n}\rangle$.
Dengan menggunakan fakta bahwa $\lambda\approx\frac{\left|p_{n+1}-p\right|}{\left|p_{n}-p\right|^{\alpha}}$,
kita peroleh $\lambda=0.8$ untuk setiap barisan. Oleh karena itu,
pada $\langle q_{n}\rangle$, setiap suku seharusnya memiliki $-\log0.8\approx.1$
angka signifikan akurat lebih banyak daripada suku sebelumnya. Dalam
istilah yang lebih mudah dipahami, barisan itu memperoleh sekitar
satu angka signifikan akurat tambahan setiap sepuluh suku. Hal ini
tampak dari pengamatan bahwa $q_{0}$ memiliki galat sekitar $3(10)^{-1}$,
sedangkan $q_{10}$ memiliki galat sekitar $3(10)^{-2}$. Untuk $\langle r_{n}\rangle$,
setiap suku diperkirakan memiliki sekitar $-\log(0.8\cdot e^{.322})\approx-0.04$
angka signifikan akurat lebih banyak daripada $1.322$ kali banyaknya
angka signifikan akurat pada suku sebelumnya. Sebagai contoh, $r_{3}$
memiliki sekitar $\log\left|\frac{e}{2.145(10)^{-2}}\right|\approx2.1$
angka signifikan akurat, sedangkan $r_{4}$ memiliki sekitar $1.322(2.1)-.04\approx2.73$
angka signifikan akurat, $r_{5}$ memiliki $1.322(2.73)-.04\approx3.57$
angka signifikan akurat, dan seterusnya hingga $r_{8}$ memiliki sekitar
$8.1$ angka signifikan akurat. Hal ini kembali ditegaskan oleh tabel,
karena $\log\left|\frac{e}{r_{8}-e}\right|=\log\left|\frac{e}{2.113(10)^{-8}}\right|\approx8.1$.
Walaupun kita dapat melakukan perhitungan serupa untuk $\langle t_{n}\rangle$,
lebih mudah untuk menilainya secara langsung karena yang perlu kita
perhatikan hanyalah bahwa eksponen dalam notasi ilmiah kira-kira menjadi
dua kali lipat dari satu suku ke suku berikutnya. Memang demikian:
eksponennya berubah dari 1 menjadi 2, lalu 3, 6, 11, dan seterusnya. 

Perhatikan bahwa dalam seluruh analisis ini, kita telah mengabaikan
syarat bahwa $n$ harus ``besar''. Hal itu dapat diterima dalam
kasus ini karena barisan-barisan tersebut sengaja dirancang sedemikian
rupa sehingga bahkan $n=0$ pun sudah cukup besar! Dalam penerapan
praktis, tidak demikian.

Untuk memahami betapa jauh lebih cepat suatu orde kekonvergenan dibanding
orde lainnya, tinjau kembali relasi 
\[
d(p_{n+1})\approx\alpha d(p_{n})-\log\left(\lambda|p|^{\alpha-1}\right)
\]
. Sekarang, anggaplah kita mengetahui bahwa $d(p_{n_{0}})=d_{n_{0}}$
untuk suatu $n_{0}$ tertentu yang cukup besar sehingga hampiran tersebut
masuk akal. Maka dapat ditunjukkan bahwa, untuk $\alpha>1$, 
\[
d(p_{n_{0}+k})\approx(d_{n_{0}}-C)\alpha^{k}+C
\]
dengan ${\displaystyle C=\frac{\log\left(\lambda|p|^{\alpha-1}\right)}{\alpha-1}}$. 

\begin{digression}{Menyelesaikan Relasi Rekurensi} 

\noindent Relasi $d(p_{n+1})\approx\alpha d(p_{n})-\log\left(\lambda|p|^{\alpha-1}\right)$
merupakan contoh relasi rekurensi. Secara khusus, relasi tersebut
merupakan relasi rekurensi linear takhomogen orde pertama dengan koefisien
konstan karena berbentuk
\[
a_{n+1}=k_{1}a_{n}+k_{2}
\]
dengan $k_{1}$ dan $k_{2}$ adalah konstanta. Relasi rekurensi linear
takhomogen dapat diselesaikan dengan menjumlahkan solusi homogen dan
solusi partikular. Untuk solusi partikular, kita mencari solusi berbentuk
$a_{n}=A$ (untuk semua $n$) dengan menyubstitusikan solusi yang
diasumsikan ini ke dalam relasi rekurensi. Penyubstitusian itu menghasilkan
$A=k_{1}A+k_{2}$, sehingga $A=\frac{k_{2}}{1-k_{1}}$ merupakan solusi
semacam itu. Untuk solusi homogen, kita mencari barisan berbentuk
$a_{n}=r^{n}$ yang memenuhi $a_{n+1}=k_{1}a_{n}+0$. Dengan menyubstitusikan
solusi yang diasumsikan ke dalam relasi rekurensi yang telah diubah
(homogen), kita peroleh $r^{n+1}=k_{1}r^{n}$. Dengan menata ulang,
$r^{n}(r-k_{1})=0$ sehingga $r=0$ atau $r=k_{1}$. Perhatikan bahwa
$Bk_{1}^{n}$ juga merupakan solusi untuk sembarang konstanta $B$.
Hal ini mencakup solusi $a_{n}=0$ yang diperoleh dengan menetapkan
$r=0$. Akhirnya, dengan menggabungkan solusi partikular dan solusi
homogen, solusi dari $a_{n+1}=k_{1}a_{n}+k_{2}$ adalah $a_{n}=Bk_{1}^{n}+\frac{k_{2}}{1-k_{1}}$
untuk sembarang konstanta $B$. Dalam kasus $d(p_{n+1})\approx\alpha d(p_{n})-\log\left(\lambda|p|^{\alpha-1}\right)$,
$k_{1}=\alpha$ dan $k_{2}=-\log\left(\lambda|p|^{\alpha-1}\right)$
sehingga $d(p_{n})=B\alpha^{n}+\frac{\log\left(\lambda|p|^{\alpha-1}\right)}{\alpha-1}$.
Nilai $B$ ditentukan dengan menyubstitusikan sembarang unsur barisan
yang diketahui ke dalam rumus ini, lalu menyelesaikannya untuk $B$.
Jika $d(p_{n_{0}})=d_{n_{0}}$, diperoleh $d(p_{n})=\left(d_{n_{0}}-\frac{\log\left(\lambda|p|^{\alpha-1}\right)}{\alpha-1}\right)\alpha^{n}+\frac{\log\left(\lambda|p|^{\alpha-1}\right)}{\alpha-1}$. 

\noindent\end{digression}  

\noindent Hal penting yang perlu diperhatikan di sini ialah bahwa
$d(p_{n_{0}+k})$ merupakan fungsi eksponensial ketika $\alpha>1$.
Banyaknya angka signifikan yang akurat bertumbuh secara eksponensial
dengan basis $\alpha$. Seperti yang telah kita lihat, untuk $\alpha=1$,
banyaknya angka signifikan bertumbuh secara linear. Dalam kalkulus,
Anda telah mempelajari bahwa setiap fungsi eksponensial bertumbuh
jauh lebih cepat daripada setiap fungsi polinomial. Oleh karena itu,
wajar dan benar untuk menyimpulkan bahwa barisan yang konvergen dengan
orde lebih besar daripada $1$ berkonvergensi jauh lebih cepat daripada
barisan dengan orde linear ($\alpha=1$). 

Namun, berhati-hatilah. Berdasarkan pengetahuan kalkulus yang sama,
Anda juga akan menyimpulkan bahwa barisan $\langle2^{-n}\rangle$
konvergen ke 0 jauh lebih cepat daripada $\langle n^{-2}\rangle$.
Menurut sebagian ukuran, hal itu benar, tetapi tidak menurut semua
ukuran. Tinjau orde kekonvergenan kedua barisan tersebut. Kita mencari
nilai $\alpha_{1}$ dan $\alpha_{2}$ sedemikian sehingga
\[
\lim_{n\to\infty}\frac{|2^{-(n+1)}-0|}{|2^{-n}-0|^{\alpha_{1}}}=\lambda_{1}\qquad\mbox{dan}\qquad\lim_{n\to\infty}\frac{|(n+1)^{-2}-0|}{|n^{-2}-0|^{\alpha_{2}}}=\lambda_{2}
\]
untuk suatu bilangan real $\lambda_{1}$ dan $\lambda_{2}$. Sedikit
aljabar menghasilkan bentuk berikut:
\[
\begin{gathered}\frac{|2^{-(n+1)}-0|}{|2^{-n}-0|^{\alpha_{1}}}=\frac{2^{-n-1}}{2^{-\alpha_{1}n}}=2^{(\alpha_{1}-1)n-1}\\
\mbox{while}\\
\frac{|(n+1)^{-2}-0|}{|n^{-2}-0|^{\alpha_{2}}}=\frac{n^{2\alpha_{2}}}{n^{2}+2n+1}.
\end{gathered}
\]
Satu-satunya cara agar $\lim_{n\to\infty}2^{(\alpha_{1}-1)n-1}$ menjadi
konstanta taknol ialah jika $\alpha_{1}=1$. Satu-satunya cara agar
$\lim_{n\to\infty}\frac{n^{2\alpha_{2}}}{n^{2}+2n+1}$ menjadi konstanta
taknol ialah jika koefisien utama pembilang dan penyebut sama. Artinya,
$\alpha_{2}$ juga harus bernilai 1. Jadi, $\langle2^{-n}\rangle$
dan $\langle n^{-2}\rangle$ keduanya konvergen ke nol dengan orde
linear. Menurut ukuran ini, keduanya sama-sama sangat lambat berkonvergensi!
Namun, terasa ada yang kurang tepat jika kita menyatakan bahwa $\langle2^{-n}\rangle$
dan $\langle n^{-2}\rangle$ berkonvergensi dengan kecepatan yang
sama. 

\begin{digression}{Orde Kekonvergenan yang Tidak Terdefinisi}

Tidak semua barisan konvergen memiliki orde kekonvergenan yang terdefinisi.
Barisan ${\displaystyle \left\langle \frac{e^{n}}{e^{e^{n}}}\right\rangle }$,
yang konvergen ke 0, merupakan salah satu contohnya.
\begin{align*}
\frac{\left|\frac{e^{n+1}}{e^{e^{n+1}}}-0\right|}{\left|\frac{e^{n}}{e^{e^{n}}}-0\right|^{\alpha}} & =\frac{\left|\frac{e^{n+1}}{e^{e^{n+1}}}-0\right|}{\left|\frac{e^{n}}{e^{e^{n}}}-0\right|^{\alpha}}\\
 & =\frac{\left|e^{n+1-e^{n+1}}\right|}{\left|e^{\alpha n-\alpha e^{n}}\right|}\\
 & =\left|e^{1+(1-\alpha)n+(\alpha-e)e^{n}}\right|
\end{align*}
Jika kita menetapkan $\alpha\leq e$, maka ${\displaystyle \lim_{n\to\infty}[1+(1-\alpha)n+(\alpha-e)e^{n}]=-\infty}$,
yang mengakibatkan ${\displaystyle \lim_{n\to\infty}\left|e^{1+(1-\alpha)n+(\alpha-e)e^{n}}\right|=0}$.
Jika kita menetapkan $\alpha>e$, maka ${\displaystyle \lim_{n\to\infty}[1+(1-\alpha)n+(\alpha-e)e^{n}]=\infty}$,
yang mengakibatkan ${\displaystyle \lim_{n\to\infty}\left|e^{1+(1-\alpha)n+(\alpha-e)e^{n}}\right|=\infty}$.
Tidak ada nilai $\alpha$ yang membuat limit ini menjadi konstanta
taknol.

\end{digression}

\subsubsection{Laju kekonvergenan suatu barisan\index{kekonvergenan!laju}}

Untuk barisan yang berkonvergensi dengan orde linear, kita memerlukan
ukuran yang lebih halus daripada orde untuk menentukan mana yang lebih
cepat. Ingat kembali dari kalkulus, 
\begin{eqnarray*}
\lim_{n\to\infty}\frac{2^{-n}}{n^{-2}} & = & \lim_{n\to\infty}\frac{n^{2}}{2^{n}}\\
 & = & \lim_{n\to\infty}\frac{2n}{2^{n}\ln2}\\
 & = & \lim_{n\to\infty}\frac{2}{2^{n}(\ln2)^{2}}=0,
\end{eqnarray*}
yang menunjukkan bahwa $\langle2^{-n}\rangle$ mendekati 0 jauh lebih
cepat daripada $\langle n^{-2}\rangle$. Anda mungkin juga mengingat
perbandingan antara fungsi pangkat:
\[
\lim_{n\to\infty}\frac{n^{-p}}{n^{-q}}=0
\]
jika $p>q>0$; antara fungsi eksponensial:
\[
\lim_{n\to\infty}\frac{a^{-n}}{b^{-n}}=0
\]
jika $a>b\geq1$; dan antara keduanya:
\[
\lim_{n\to\infty}\frac{a^{-n}}{n^{-q}}=0
\]
jika $a>1$. Dengan kata lain, barisan berbentuk $\left\langle \frac{1}{a^{n}}\right\rangle $
berkonvergensi ke nol lebih cepat daripada barisan berbentuk $\left\langle \frac{1}{n^{p}}\right\rangle $
jika $a>1$. Barisan $\left\langle \frac{1}{a^{n}}\right\rangle $
berkonvergensi ke nol lebih cepat daripada $\left\langle \frac{1}{b^{n}}\right\rangle $
jika $a>b\geq1$. Barisan $\left\langle \frac{1}{n^{p}}\right\rangle $
berkonvergensi ke nol lebih cepat daripada $\left\langle \frac{1}{n^{q}}\right\rangle $
jika $p>q>0$. Tidak semua fungsi sesederhana ini, tetapi kita dapat
menggunakannya sebagai tolok ukur. Misalkan $\langle p_{n}\rangle$
berkonvergensi ke $p$, $\langle b_{n}\rangle$ berkonvergensi ke
0 dan $|p_{n}-p|\leq\lambda|b_{n}|$ untuk suatu konstanta $\lambda$
dan semua $n$ yang cukup besar. Maka kita katakan bahwa $\langle p_{n}\rangle$
berkonvergensi ke $p$ dengan laju kekonvergenan \index{kekonvergenan!laju}
$O(b_{n})$, dibaca ``O besar dari $b_{n}$''\index{O besar}. Karena
kita sudah mengenal barisan berbentuk $\left\langle \frac{1}{a^{n}}\right\rangle $
untuk suatu konstanta $a>1$ dan $\left\langle \frac{1}{n^{p}}\right\rangle $
untuk suatu konstanta $p>0$, dan keduanya cukup sederhana, biasanya
$\langle b_{n}\rangle$ akan berupa salah satunya. Sebagai contoh,
$\left\langle \frac{2n+1}{4n}\right\rangle $ berkonvergensi ke $\frac{1}{2}$,
dan
\[
\left|\frac{2n+1}{4n}-\frac{1}{2}\right|=\frac{1}{4n}\leq\frac{1}{4}\cdot\frac{1}{n},
\]
sehingga $\left\langle \frac{2n+1}{4n}\right\rangle $ berkonvergensi
dengan laju $O(\frac{1}{n})$. Kita juga dapat menuliskan $\frac{2n+1}{4n}=\frac{1}{2}+O(\frac{1}{n})$
untuk menyampaikan hal yang persis sama. Biasanya, ketika mencari
laju kekonvergenan, kita berusaha menemukan barisan dengan kekonvergenan
tercepat dari kumpulan contoh sederhana kita yang memenuhi definisi
tersebut. Dalam kasus ini, tidak ada yang lebih cepat.

Pada dasarnya, semua barisan yang dipelajari secara mendalam dalam
kalkulus konvergen dengan orde linear. Jadi, apa yang diperlukan agar
suatu barisan berkonvergensi dengan orde yang lebih tinggi? Mari kita
tinjau $\left\langle e^{-2^{n}}\right\rangle $. 
\[
\lim_{n\to\infty}\frac{|e^{-2^{n+1}}-0|}{|e^{-2^{n}}-0|^{\alpha}}=\lim_{n\to\infty}\frac{e^{-2\cdot2^{n}}}{e^{-\alpha2^{n}}}=1
\]
ketika $\alpha=2$. Jadi, $\left\langle \frac{1}{e^{2^{n}}}\right\rangle $
berkonvergensi secara kuadratik. Pada dasarnya, diperlukan eksponen
yang bertumbuh secara eksponensial agar suatu barisan berkonvergensi
dengan orde lebih besar daripada $1$.

\begin{digression}{Menghampiri $\pi$} \index{pi@$\pi$!hampiran}

\noindent Barisan 
\[
\frac{1103\cdot2^{3/2}}{9801},\ \frac{1130173253125}{313826716467\cdot2^{7/2}},\ \frac{1029347477390786609545}{1116521080257783321\cdot2^{23/2}},\ldots
\]
konvergen ke $\frac{1}{\pi}$. Suku-sukunya diberikan oleh rumus 
\[
\left\langle \frac{\sqrt{8}}{9801}\sum_{j=0}^{n}\frac{(4j)!(1103+26390j)}{(j!)^{4}\cdot396^{4j}}\right\rangle _{n=0,1,2,3,\ldots}
\]
karya Srinivasa Ramanujan\index{Ramanujan!Srinivasa}. Untuk keperluan
praktis, barisan ini berkonvergensi sangat cepat. Suku pertamanya
sudah memiliki sekitar 8 angka signifikan yang akurat:
\begin{eqnarray*}
\frac{1103\cdot2^{3/2}}{9801} & \approx & 0.31830987844047012321768445317\\
\frac{1}{\pi} & \approx & 0.31830988618379067153776752674,
\end{eqnarray*}
sedangkan suku keduanya memiliki sekitar 16:
\[
\left|\frac{1130173253125}{313826716467\cdot2^{7/2}}-\frac{1}{\pi}\right|\approx6.48(10)^{-17},
\]
dua kali akurasi suku pertama. Akurasi suku ketiganya bahkan sudah
melampaui presisi ganda. 

Sangat menggoda untuk percaya, atau berharap, bahwa barisan ini berkonvergensi
secara kuadratik, tetapi ternyata tidak. Suku ketiganya akurat hingga
sekitar 24 angka signifikan. Setiap suku dalam barisan kira-kira 8
angka signifikan lebih akurat daripada suku sebelumnya---ciri khas
barisan yang berkonvergensi secara linear. 

\end{digression}

\subsection{Konsep Utama}
\begin{description}
\item [{Orde~kekonvergenan~barisan\index{kekonvergenan!orde}:}] Barisan
$\langle p_{n}\rangle$ konvergen ke $p$ dengan orde kekonvergenan
$\alpha\geq1$ jika
\[
\lim_{n\to\infty}\frac{\left|p_{n+1}-p\right|}{\left|p_{n}-p\right|^{\alpha}}=\lambda
\]
untuk suatu bilangan real $\lambda>0$. 
\item [{Galat~absolut:}] Untuk barisan $\langle p_{n}\rangle$ yang konvergen
ke $p$ dengan orde $\alpha$, galat absolut suku-suku berurutan dihubungkan
oleh hampiran
\[
|p_{n+1}-p|\approx\lambda|p_{n}-p|^{\alpha}
\]
untuk $n$ yang cukup besar.
\item [{Angka~signifikan~yang~akurat:}] \index{akurasi!angka signifikan}Untuk
barisan $\langle p_{n}\rangle$ yang konvergen ke $p$ dengan orde
$\alpha$, banyaknya angka signifikan yang akurat pada suku-suku berurutan
dihubungkan oleh hampiran
\[
d(p_{n+1})\approx\alpha d(p_{n})-\log\left(\lambda|p|^{\alpha-1}\right)
\]
untuk $n$ yang cukup besar. Dalam bentuk tertutup (untuk $\alpha\neq1$)
\[
d(p_{n+k})=(d_{n}-C)\alpha^{k}+C
\]
dengan ${\displaystyle C=\frac{\log\left(\lambda|p|^{\alpha-1}\right)}{\alpha-1}}$.
\item [{Laju~kekonvergenan~barisan\index{O besar}\index{kekonvergenan!laju}:}] Barisan
$\langle p_{n}\rangle$ konvergen ke $p$ dengan laju kekonvergenan
$O(b_{n})$ jika $\langle b_{n}\rangle$ konvergen ke 0 dan 
\[
|p_{n}-p|\leq\lambda|b_{n}|
\]
untuk suatu konstanta $\lambda$ dan semua $n$ yang cukup besar.
\end{description}

\subsection{Octave}

Perulangan merupakan alat yang sangat berguna dalam segala jenis pemrograman.
Ketika Anda perlu menjalankan suatu prosedur berulang kali dengan
masukan yang berbeda-beda, perulangan mungkin merupakan solusi yang
tepat. Meskipun Octave menyediakan beberapa jenis perulangan, untuk
saat ini kita hanya akan membahas perulangan \texttt{for} \index{Octave!perulangan for@perulangan \texttt{for}}.
Gagasannya adalah menggunakan suatu variabel, yang kadang-kadang disebut
pencacah, untuk menghitung berapa kali prosedur telah dijalankan.
Ketika prosedur telah dijalankan sebanyak yang diinginkan, perulangan
berakhir dan program berlanjut dari sana. Hampir pasti Anda sudah
menjumpai gagasan ini bahkan sebelum pernah menulis program komputer.
Jika Anda pernah pergi ke pasar malam dan membayar satu dolar untuk
melempar selusin gelang dengan harapan salah satunya melingkari leher
botol soda, Anda telah mengalami perulangan. Anda mungkin bahkan menghitung
gelang-gelang itu saat melemparkannya. Andalah pencacahnya! Anda harus
menjalankan prosedur melemparkan gelang ke deretan botol sebanyak
12 kali. Jadi, mungkin Anda melempar gelang pertama dan menghitung
dalam hati ``1''. Kemudian Anda melempar gelang berikutnya dan menghitung
``2''. Lalu gelang berikutnya lagi dan menghitung ``3''. Begitu
seterusnya hingga ``12''. Setelah gelang terakhir dilempar, Anda
melanjutkan kegiatan di pasar malam. 

Perulangan \texttt{for} merupakan analogi abstrak dari situasi ini.
Misalkan Anda ingin menghitung $1!$, $2!$, $3!$, dan seterusnya
hingga $12!$. Di Octave, Anda dapat membuat berkas \texttt{.m} berikut
lalu menjalankannya. \begin{verbatim}
factorial(1)
factorial(2)
factorial(3)
factorial(4)
factorial(5)
factorial(6)
factorial(7)
factorial(8)
factorial(9)
factorial(10)
factorial(11)
factorial(12)
\end{verbatim}

\noindent Namun, cara ini dapat menjemukan dan kurang ramah bagi pembaca,
khususnya jika kita ingin melakukan suatu perhitungan jauh lebih dari
12 kali. Tujuan perulangan adalah mengurangi pengulangan dalam pendekatan
ini. Kita ingin menjalankan prosedur untuk menghitung faktorial 12
bilangan bulat yang berbeda, sehingga perulangan cocok digunakan.
Sintaks perulangan ialah menyiapkan pencacah, menulis kode untuk menjalankan
prosedur, lalu menandai akhir \index{Octave!end@\texttt{end}} perulangan.
Bentuknya kira-kira seperti berikut. \begin{verbatim}
for j=first:last
  do something. 
end%for
\end{verbatim}

\noindent Perintah ini akan menyebabkan Octave menjalankan prosedur
sekali untuk setiap bilangan bulat dari \texttt{first} hingga \texttt{last},
termasuk \texttt{first} dan \texttt{last}. Nilai pencacah, dalam hal
ini \texttt{j}, dapat digunakan dalam prosedur. Jadi, untuk menghitung
$1!$ hingga $12!$, kita dapat menulis \begin{verbatim}
for j=1:12
  factorial(j)
end%for
\end{verbatim}

\noindent Program ini akan menghasilkan keluaran yang persis sama
dengan program yang memiliki satu baris untuk setiap faktorial. Jika
kelak Anda ingin menghitung $1!$ hingga $20!$ sebagai gantinya,
Anda hanya perlu mengubah 12 menjadi 20. Perulangan \texttt{for} adalah
sahabat Anda!

Sekarang misalkan kita ingin menghitung $\alpha$ untuk setiap kelompok
tiga nilai $|s_{n}-e|$ yang berurutan pada Tabel \ref{tab:absoluteErrors}.
Karena ada 9 kelompok seperti itu, kita perlu membuat perulangan yang
dijalankan 9 kali. Di dalam perulangan tersebut, kita perlu melakukan
perhitungan ${\displaystyle \alpha=\frac{\ln\left|\frac{s_{n+2}-e}{s_{n+1}-e}\right|}{\ln\left|\frac{s_{n+1}-e}{s_{n}-e}\right|}}$.
Namun, sebelum memulai, kita perlu memberi tahu Octave tentang 11
nilai dari tabel tersebut. Cara paling praktis ialah menyimpannya
dalam sebuah larik. Larik \index{Octave!larik} menyerupai vektor
dan memiliki komponen. Dalam hal ini, setiap komponen akan menyimpan
satu nilai dari tabel. Sintaks untuk membuat larik juga sangat mirip
dengan notasi vektor. Kita akan menggunakan kurung siku untuk mengapit
komponen-komponen larik dan memisahkannya dengan koma. Jadi, baris
pertama kode Octave kita akan tampak seperti berikut. \begin{verbatim}
errs = [2.817*10^(-1), 1.03*10^(-1), 2.022*10^(-2), 1.451*10^(-3), ...
        2.046*10^(-5), 2.07*10^(-8), 2.953*10^(-13), 4.263*10^(-21), ...
        8.777*10^(-34), 2.608*10^(-54), 1.595*10^(-87)]
\end{verbatim}Tanda elipsis (tiga titik berturut-turut) di akhir dua baris pertama
diperlukan untuk memberi tahu Octave bahwa perintah tersebut berlanjut
ke baris berikutnya. Tanpa tanda itu, memecah satu perintah menjadi
beberapa baris akan menimbulkan galat sintaks. Di Octave, memulai
baris baru juga merupakan isyarat untuk memulai perintah baru.

Sekarang Octave mengetahui nilai-nilai $|s_{n}-e|$. Menggunakan vektor
ini sangat mirip dengan menggunakan subskrip. Nilai pertama, $2.817(10)^{-1}$,
disebut \texttt{errs(1)}. Nilai kedua disebut \texttt{errs(2)}. Nilai
ketiga disebut \texttt{errs(3),} dan seterusnya. Panjang larik \texttt{errs}
dapat diperoleh dengan fungsi \texttt{length()} \index{Octave!panjang larik}
di Octave. Perintah \texttt{length(errs)-2} akan digunakan alih-alih
menuliskan angka 9 langsung di dalam kode. Jadi, kita dapat menyelesaikan
kode Octave sebagai berikut. \begin{verbatim}
errs = [2.817*10^(-1), 1.03*10^(-1), 2.022*10^(-2), 1.451*10^(-3), ...
        2.046*10^(-5), 2.07*10^(-8), 2.953*10^(-13), 4.263*10^(-21), ...
        8.777*10^(-34), 2.608*10^(-54), 1.595*10^(-87)];
for j=1:length(errs)-2
  alpha = log(errs(j+2)/errs(j+1))/log(errs(j+1)/errs(j))
end%for
\end{verbatim}

\noindent Kode ini menghasilkan keluaran berikut: \begin{verbatim}
alpha =  1.6182
alpha =  1.6181
alpha =  1.6176
alpha =  1.6182
alpha =  1.6180
alpha =  1.6180
alpha =  1.6180
alpha =  1.6180
alpha =  1.6180
\end{verbatim}Lumayan, tetapi kita dapat membuatnya lebih baik. Mari kita hitung
$\alpha$, $\lambda$, dan $d(s_{n})$ dengan dua metode berbeda---secara
langsung dan menggunakan rumus $d(p_{n+1})\approx\alpha d(p_{n})-\log\left(\lambda|p|^{\alpha-1}\right)$.
Kemudian, mari tampilkan hasilnya dalam tabel yang diformat dengan
rapi. 

Kita memerlukan perintah \texttt{disp()} dan larik berindeks dua.
Perintah \texttt{disp()} \index{Octave!disp@\texttt{disp}} digunakan
untuk menampilkan teks atau suatu besaran. Jika digunakan untuk teks,
teks harus diapit tanda petik tunggal; jika digunakan untuk besaran,
tanda petik itu tidak diperlukan. Jadi, program Octave dapat menampilkan
kata ``hello'' dengan perintah \texttt{disp('hello')} atau menampilkan
nilai $\ln(2)$ dengan perintah \texttt{disp(log(2))}. Perintah \texttt{disp()}
juga dapat menangani variabel. Jadi, jika \texttt{p1} dan \texttt{p2}
telah diberi nilai, kita dapat menampilkan selisihnya dengan \texttt{disp(p2-p1)}.
Larik berindeks dua dapat dipandang sebagai tabel atau matriks. Larik
ini menyimpan nilai dalam susunan baris dan kolom. Jadi, alih-alih
\texttt{errs(j)} seperti sebelumnya, kita dapat memiliki \texttt{errs(j,k)},
dengan \texttt{j} menyatakan baris dan \texttt{k} menyatakan kolom.
Program \begin{verbatim}
A(2,4) = 7;
disp(A);
\end{verbatim}menghasilkan \begin{verbatim}
   0   0   0   0
   0   0   0   7 
\end{verbatim}

\noindent Baiklah, mari kembali ke persoalan yang sedang kita kerjakan.
Kita akan menggabungkan semua yang telah kita pelajari tentang Octave
dalam satu program. \begin{verbatim}
errs = [2.817*10^(-1), 1.03*10^(-1), 2.022*10^(-2), 1.451*10^(-3), ...
     2.046*10^(-5), 2.07*10^(-8), 2.953*10^(-13), 4.263*10^(-21), ...
     8.777*10^(-34), 2.608*10^(-54), 1.595*10^(-87)];

d = @(x) -log(x/exp(1))/log(10);
for j=1:9
    % alpha:
  T(j,1) = log(errs(j+2)/errs(j+1))/log(errs(j+1)/errs(j));
    % lambda:
  T(j,2) = errs(j+2)/errs(j+1)^T(j,1);
    % d (explicit):
  T(j,3) = d(errs(j+2));
end

alpha = 1.61804;
lambda = 0.8;
constant = log(lambda*exp(alpha-1))/log(10);
T(1,4) = T(1,3);
for j=2:9
    % d (recursive)
  T(j,4) = alpha * T(j-1,4) - constant;
end%for

disp('      alpha     lambda   d (expl)    d (rec)');
disp('  --------------------------------------------');
disp(T) 
\end{verbatim} menghasilkan \begin{verbatim}
      alpha     lambda   d (expl)    d (rec)
  --------------------------------------------
    1.61816    0.80015    2.12851    2.12851
    1.61814    0.80010    3.27263    3.27252
    1.61764    0.79855    5.12339    5.12356
    1.61822    0.80158    8.11832    8.11863
    1.61797    0.79941   12.96403   12.96477
    1.61804    0.80045   20.80458   20.80601
    1.61805    0.80059   33.49095   33.49346
    1.61804    0.80031   54.01799   54.02225
    1.61804    0.80031   87.23153   87.23866 
\end{verbatim}Ada baiknya meluangkan waktu untuk memastikan bahwa Anda memahami
semua baris dalam program ini. Program ini menggunakan penugasan,
fungsi bawaan, fungsi anonim, keluaran sederhana, larik, dan perulangan
\texttt{for}. Tanda \texttt{\%} \index{Octave!%@\texttt{\%}} pada
suatu baris memberi tahu Octave agar mengabaikan tanda itu serta seluruh
isi setelahnya pada baris yang sama. Bagian tersebut disebut komentar\index{Octave!komentar}.
Komentar ditujukan semata-mata bagi pembaca manusia untuk mendokumentasikan
apa yang dilakukan program. Program yang panjang harus selalu didokumentasikan
agar setiap penggunanya lebih mudah memahami apa yang dilakukan program
tersebut. Di sini komentarnya sederhana, tetapi komentar dapat jauh
lebih terperinci.

\begin{digression}{Lebih lanjut tentang perulangan \texttt{for}}

\noindent Sintaks \texttt{first:last} yang digunakan dalam perulangan
\texttt{for} sebenarnya merupakan bentuk ringkas untuk larik yang
berisi bilangan bulat dari \texttt{first} hingga \texttt{last}. Bilangan
ketiga dapat ditambahkan ke notasi ini untuk menghasilkan barisan
aritmetika sebarang seperti $10,13,16,19,22$ dan $1,1.1,1.2,1.3,1.4,1.5,1.6,1.7,1.8,1.9,2$.
Bentuk lengkap notasi ini adalah \texttt{first:step:last}; jadi, \texttt{10:3:22}
akan menghasilkan larik yang berisi suku-suku $10,13,16,19,22$. Selain
itu, sintaks larik ini sepenuhnya tidak bergantung pada penggunaannya
dalam perulangan \texttt{for}. Larik dapat dibentuk dengan sintaks
ini untuk tujuan apa pun, seperti diperlihatkan di sini.\begin{verbatim}
octave:1> a=1:.1:2
a =
 Columns 1 through 7:
    1.0000    1.1000    1.2000    1.3000    1.4000    1.5000    1.6000
 Columns 8 through 11:
    1.7000    1.8000    1.9000    2.0000
octave:2> b=[1,1.1,1.2,1.3,1.4,1.5,1.6,1.7,1.8,1.9,2]
b =
 Columns 1 through 8:
   1.0000   1.1000   1.2000   1.3000   1.4000   1.5000   1.6000   1.7000
 Columns 9 through 11:
   1.8000   1.9000   2.0000
octave:3> a==b % Check to see if arrays a and b are equal--they are not!
ans =
  1  1  1  1  1  1  1  0  1  1  1
octave:4> format('bit') % display numbers as stored in memory
octave:5> a(8) % The eighth term of array a (1.7)
ans = 0011111111111011001100110011001100110011001100110011001100110100
octave:6> b(8) % The eighth term of array b (1.7)
ans = 0011111111111011001100110011001100110011001100110011001100110011
\end{verbatim}

\noindent Menarik untuk diperhatikan bahwa bilangan $1.7$ yang dihasilkan
dengan notasi barisan aritmetika (seperti pada variabel \texttt{a})
berbeda dari bilangan $1.7$ yang dimasukkan dengan menuliskannya
langsung sebagai konstanta (seperti pada variabel \texttt{b}). Kedua
representasi $1.7$ tersebut berbeda pada tiga digit biner terakhirnya!

Kembali ke perulangan \texttt{for}, cara larik tersebut dibuat tidaklah
penting. Perulangan \texttt{for} akan menjalankan kodenya sekali untuk
setiap nilai dalam larik. Sebagai contoh, program\begin{verbatim}
for i=[1,5,7,0,8]
  disp(factorial(i))
end%for
\end{verbatim}menghasilkan\begin{verbatim}
 1
 120
 5040
 1
 40320
\end{verbatim}

\end{digression}

\noindent\startexercises

\subsection{Latihan}
\begin{enumerate}
\item Berikut diberikan beberapa barisan konvergen beserta limitnya. Tentukan
orde kekonvergenan masing-masing.

\begin{enumerate}
\item ${\displaystyle \left\langle \frac{n!}{n^{n}}\right\rangle \to0}$
\item \label{exc:convergence}${\displaystyle \left\langle \frac{1}{3^{e^{n}}}\right\rangle \to0}$ \hasasolution 
\item \label{exc:convergence-1}${\displaystyle \left\langle \frac{2^{2^{n}}-2}{2^{2^{n}}+3}\right\rangle \to1}$ \hasasolution 
\item \label{exc:convergence-2}${\displaystyle \left\langle \frac{n^{2}}{1+n^{2}}\right\rangle \to1}$ \hasananswer 
\item ${\displaystyle \left\langle f^{\circ n}(1)\right\rangle \to\sqrt{3}}$
dengan ${\displaystyle f(x)=\frac{x^{2}+3}{2x}}$. Catatan: $f^{\circ n}$
berarti $f$ dikomposisikan dengan dirinya sendiri sebanyak $n$ kali.
\end{enumerate}
\item \label{enu:linearConvergenceExample}Tunjukkan bahwa barisan ${\displaystyle \left\langle \frac{n+1}{n-1}\right\rangle }$
konvergen ke 1 secara linear.
\item Tunjukkan bahwa barisan $p_{n}=2^{1-2^{n}}$ berkonvergensi secara
kuadratik.
\item Berikan contoh barisan yang konvergen ke 0 dengan orde $\alpha=10$.
\item Perkirakan orde kekonvergenan barisan $p_{n}$ dan jelaskan jawaban
Anda.

\begin{center}
\begin{tabular}{c|c|c|c}
$n$ & $\frac{|p_{n+1}-p|}{|p_{n}-p|^{1.2}}$ & $\frac{|p_{n+1}-p|}{|p_{n}-p|^{1.3}}$ & $\frac{|p_{n+1}-p|}{|p_{n}-p|^{1.4}}$\tabularnewline
\hline 
25 & $9.07(10)^{-6}$ & $.0110$ & $13.39$\tabularnewline
26 & $1.88(10)^{-7}$ & $.00303$ & $48.65$\tabularnewline
27 & $1.01(10)^{-9}$ & $.000530$ & $277.8$\tabularnewline
28 &  &  & \tabularnewline
\end{tabular}
\par\end{center}
\item Berikut diberikan beberapa barisan yang konvergen secara linear beserta
limitnya. Untuk masing-masing, tentukan laju kekonvergenan (tercepat)
yang berbentuk $O\left(\frac{1}{n^{p}}\right)$ atau $O\left(\frac{1}{a^{n}}\right)$.
Jika tidak mungkin, usulkan laju kekonvergenan yang masuk akal.

\begin{enumerate}
\item ${\displaystyle 6,\frac{6}{7},\frac{6}{49},\frac{6}{343},\frac{6}{2401},\ldots\to0}$
\item ${\displaystyle \left\langle \frac{11n-2}{n+3}\right\rangle \to11}$
\item \label{exc:convergence-3}${\displaystyle \left\langle \frac{\sin n}{\sqrt{n}}\right\rangle \to0}$ \hasasolution 
\item \label{exc:convergence-4}${\displaystyle \left\langle \frac{4}{10^{n}+35n+9}\right\rangle \to0}$ \hasasolution 
\item ${\displaystyle \left\langle \frac{1}{n}+\frac{1}{n^{2}}\right\rangle \to0}$
\item \label{exc:convergence-5}${\displaystyle \left\langle \frac{4}{10^{n}-35n-9}\right\rangle \to0}$ \hasasolution 
\item \label{exc:convergence-6}${\displaystyle \left\langle \frac{2n}{\sqrt{n^{2}+3n}}\right\rangle \to2}$ \hasananswer 
\item ${\displaystyle \left\langle \frac{5^{n}-2}{5^{n}+3}\right\rangle \to1}$
\item \label{exc:convergence-7}$\left\langle \sqrt{n+47}-\sqrt{n}\right\rangle \to0$ \hasananswer 
\item ${\displaystyle \left\langle \frac{n^{2}}{3n^{2}+1}\right\rangle \to\frac{1}{3}}$
\item ${\displaystyle \left\langle \frac{\pi}{e^{n}-\pi^{n}}\right\rangle \to0}$
\item \label{exc:convergence-8}${\displaystyle \left\langle \frac{n^{2}}{2^{n}}\right\rangle \to0}$ \hasasolution 
\item ${\displaystyle \left\langle \frac{7+\cos(5n)}{n^{3}+1}\right\rangle \to0}$
\item ${\displaystyle \left\langle \frac{8n^{2}}{3n^{2}+12}+\frac{n}{3n+10}\right\rangle \to3}$
\item \label{exc:convergence-9}${\displaystyle \left\langle \frac{2n^{2}+3n}{1-n^{2}}\right\rangle \to-2}$ \hasananswer 
\item ${\displaystyle \left\langle \frac{3n^{5}-5n}{1-n^{5}}\right\rangle \to-3}$
\end{enumerate}
\item Tentukan laju kekonvergenan barisan-barisan berikut saat $n\to\infty$.

\begin{enumerate}
\item ${\displaystyle \lim_{n\to\infty}\sin\frac{1}{n}=0}$
\item ${\displaystyle \lim_{n\to\infty}\sin\frac{1}{n^{2}}=0}$
\item ${\displaystyle \lim_{n\to\infty}\left(\sin\frac{1}{n}\right)^{2}=0}$
\item ${\displaystyle \lim_{n\to\infty}[\ln(n+1)-\ln(n)]=0}$
\end{enumerate}
Untuk soal \ref{exc:firstFunctionROC}-\ref{exc:lastFunctionROC},
gunakan definisi laju kekonvergenan fungsi berikut. Untuk fungsi $f(h)$,
kita mengatakan bahwa $\lim_{h\to a}f(h)=L$ memiliki laju kekonvergenan
$g(h)$ jika $|f(h)-L|\leq\lambda|g(h)|$ untuk suatu $\lambda>0$
dan semua $|h-a|$ yang cukup kecil.
\item \label{exc:firstFunctionROC} Gunakan polinomial Taylor untuk menentukan
laju kekonvergenan dari
\[
\lim_{h\to0}(2-e^{h})=1.
\]
\item Gunakan polinomial Taylor untuk menentukan laju kekonvergenan dari
\[
\lim_{h\to0}\frac{\sin(h)-e^{h}+1}{h}=0.
\]
\item Tentukan laju kekonvergenan fungsi-fungsi berikut saat $h\to0$.

\begin{enumerate}
\item ${\displaystyle \lim_{h\to0}\frac{\sin h}{h}=1}$
\item ${\displaystyle \lim_{h\to0}\frac{1-\cos h}{h}=0}$
\item ${\displaystyle \lim_{h\to0}\frac{\sin h-h\cos h}{h}=0}$
\item ${\displaystyle \lim_{h\to0}\frac{1-e^{h}}{h}=-1}$
\end{enumerate}
\item Tentukan laju kekonvergenan dari 
\[
\lim_{h\to0}\frac{h^{2}+\cos h-e^{h}}{h}=-1.
\]
\item \label{exc:lastFunctionROC}Tunjukkan bahwa 
\[
(\sin h)(1-\cos h)=0+O(h^{3}).
\]
\item \label{exc:convergence-10}\octave\ Tulislah program Octave (berkas
\texttt{.m}) yang menggunakan perulangan dan perintah \texttt{disp()}
untuk menampilkan keluaran berikut (pangkat-pangkat 7). \hasasolution 
\begin{verbatim}
 1
 7
 49
 343
 2401
 16807
 117649
 823543
 5764801
 40353607
\end{verbatim}
\item \octave\ Tulislah program Octave (berkas \texttt{.m}) yang menggunakan
perulangan dan perintah \texttt{disp()} untuk menampilkan 10 nilai
pertama dari pangkat-pangkat 5, mulai dari $5^{0}$.
\item \label{exc:convergence-11}\octave\ Tulislah program Octave (berkas
\texttt{.m}) yang menggunakan perulangan, larik, dan perintah disp()
untuk menghitung nilai ${\displaystyle f(n)=\frac{2^{2^{n}}-2}{2^{2^{n}}+3}}$
untuk $n=0,1,2,4,6,10$. \hasasolution 
\item \octave\ Tulislah program Octave (berkas \texttt{.m}) yang menggunakan
perulangan, larik, dan perintah disp() untuk menghitung nilai ${\displaystyle f(n)=\frac{2n}{\sqrt{n^{2}+3n}}}$
untuk $n=0,2,5,10,100,1000,20000$.
\item \octave\ Kode Octave berikut dimaksudkan untuk menghitung jumlah
\[
\sum_{k=1}^{30}\frac{1}{k^{2}}
\]
, tetapi tidak berhasil melakukannya. Temukan sebanyak mungkin kesalahan
pada kode tersebut. Klasifikasikan setiap kesalahan sebagai galat
kompilasi (galat yang akan mencegah program berjalan sama sekali)
atau galat logika (galat yang tidak akan mencegah program berjalan,
tetapi akan menyebabkan jumlah tersebut dihitung secara keliru).

\begin{lyxcode}
sum=1;~\\
for~k=1:30~\\
~~sum=sum+1.0/k{*}k;~\\
end~\\
disp(sum)
\end{lyxcode}
\item Beberapa barisan tidak memiliki orde kekonvergenan. Misalkan $p_{n}=\frac{2^{n}}{n!}$.

\begin{enumerate}
\item Tunjukkan bahwa $\lim_{n\to\infty}p_{n}=0$.
\item Tunjukkan bahwa $\lim_{n\to\infty}\frac{|p_{n+1}|}{|p_{n}|}=0$.
\item Tunjukkan bahwa $\left\langle \frac{|p_{n+1}|}{|p_{n}|^{\alpha}}\right\rangle $
divergen untuk setiap $\alpha>1$.
\end{enumerate}
\item Gunakan aturan praktis orde kekonvergenan untuk memperkirakan jumlah
iterasi yang diperlukan guna memperoleh hampiran bagi $\pi$ dengan
akurasi 12 angka signifikan untuk setiap orde kekonvergenan. Asumsikan
bahwa setiap barisan berawal dengan satu angka signifikan yang akurat.

\begin{enumerate}
\item $\alpha=1$, $\lambda=0.8$
\item \label{exc:convergence-12}$\alpha=1$, $\lambda=0.5$ \hasasolution 
\item $\alpha=1$, $\lambda=0.1$
\item $\alpha=1.5$
\item \label{exc:convergence-13}$\alpha=2$ \hasananswer 
\item $\alpha=3$
\end{enumerate}
\item \label{enu:orderOfConvergenceIsUnique}Buktikan bahwa orde kekonvergenan
suatu barisan bersifat unik.
\item \octave\ Tulislah perulangan \texttt{for} yang menampilkan barisan
bilangan berikut.

\begin{enumerate}
\item $7,8,9,10,11,12,13,14,15$
\item $20,19,18,17,16,15,14,13$
\item $12,12.333,12.667,13,13.333,13.667,14$
\item $1,9,25,49,81,121,169,225,289,361,441$
\item $1,.5,.25,.125,.0625,.03125,.015625$
\end{enumerate}
\end{enumerate}
\finishexercises
